% ============================================================================
%  队友章节：上游开放词汇分割   ——   负责人：史傅冠华
% ----------------------------------------------------------------------------
%  · 本文件由主文件 slides.tex 通过 % ============================================================================
%  队友章节：上游开放词汇分割   ——   负责人：史傅冠华
% ----------------------------------------------------------------------------
%  · 本文件由主文件 slides.tex 通过 % ============================================================================
%  队友章节：上游开放词汇分割   ——   负责人：史傅冠华
% ----------------------------------------------------------------------------
%  · 本文件由主文件 slides.tex 通过 % ============================================================================
%  队友章节：上游开放词汇分割   ——   负责人：史傅冠华
% ----------------------------------------------------------------------------
%  · 本文件由主文件 slides.tex 通过 \input{sections/sam.tex} 引入，
%    继承主文件导言区的全部宏 / 配色 / 字体，无需自带 \documentclass 或 \begin{document}。
%  · 内容：模型选型 → Prompt Ensemble + Rerank → 多尺度融合 + 消融 → 可视化
% ============================================================================
\section{上游分割}

% ==================== 第 1 页：模型选型与系统定位 ====================
\begin{frame}[t]{1. 上游模型选型：为什么用 SAM2？}
  \footnotesize
  \linespread{0.90}\selectfont

  \begin{columns}[T,onlytextwidth]
    \column{0.50\textwidth}
      \tcard{softorange}{2.00cm}{定位需求分析}{%
        \textbf{系统需求：}给定图像，输出 bbox / points，供 ZIM 做精细 matting。\\
        \textbf{关键约束：}
        \begin{itemize}
          \setlength{\itemsep}{0.02em}
          \item \hilite{zero-training} — 不微调模型参数；
          \item 推理高效 — 单卡秒级；
          \item bbox prompt 原生支持 — 与下游接口统一。
        \end{itemize}
        \vspace{0.02em}
        \textbf{不选 SAM3：}模型更大、推理慢约 2--3 倍，bbox-only 提升有限。\\
        \textbf{不选 Grounded-SAM-2：}额外引入 text encoder，增加显存开销。
      }
    \column{0.50\textwidth}
      \tcard{accentblue}{2.00cm}{SAM2 适配性}{%
        \begin{itemize}
          \setlength{\itemsep}{0.04em}
          \item \textbf{Hiera 架构} — 轻量高效，Hiera-Base+ 仅 80.8M 参数；
          \item \textbf{native bbox + points} — 无需适配层；
          \item \textbf{多候选 mask + IoU 评分} — 为 rerank 提供置信度；
          \item \textbf{实测性能} — A10 上约 1.16 it/s（bbox-only）；
          \item 社区活跃，checkpoint 可下载。
        \end{itemize}
      }
  \end{columns}

  \vspace{0.15em}
  \notestrip{accentblue}{\centering\scriptsize
  \textbf{定位：}SAM2 负责回答 \hilite{``目标在哪''} — 输出 bbox + points + 候选 mask；精细软 alpha 留给下游 ZIM + TTA-GF + BSP。}
\end{frame}

% ==================== 第 2 页：Prompt Ensemble + Mask Rerank ====================
\begin{frame}[t]{2. 核心优化一：Prompt Ensemble 与 IoU 引导重排序}
  \scriptsize
  \linespread{0.90}\selectfont

  \notestrip{deepred}{\tiny
  \textbf{动机：}SAM2 单次 bbox 推理易受局部纹理干扰，可能选到次优 mask。我们提出 \hilite{Prompt Ensemble + IoU 引导重排序}：多组 prompt 策略生成候选 mask，再利用 SAM2 内置 IoU 与结构先验选出最优。}

  \vspace{0.35em}
  \centering
  \resizebox{0.92\linewidth}{!}{%
  \begin{tikzpicture}[
    x=1cm, y=1cm,
    >={Latex[length=2.5mm,width=2mm]},
    every node/.style={font=\tiny},
    softsh/.style={blur shadow={shadow blur steps=5, shadow xshift=0.4pt,
                   shadow yshift=-0.5pt, shadow opacity=24, shadow blur radius=1.5pt}},
    nodebase/.style={softsh, rounded corners=3pt, line width=0.9pt, align=center, inner sep=3pt},
    io/.style={nodebase, draw=black!50, fill=black!5, minimum width=1.5cm, minimum height=0.58cm},
    sam/.style={nodebase, draw=accentblue!82!black, fill=accentblue!14, minimum width=1.55cm, minimum height=0.58cm},
    cand/.style={nodebase, draw=softorange!72!black, fill=softorange!12, minimum width=1.2cm, minimum height=0.58cm},
    ours/.style={nodebase, draw=nankai!88!black, fill=nankai!14, minimum width=1.35cm, minimum height=0.58cm},
    best/.style={nodebase, draw=softgreen!72!black, fill=softgreen!16, minimum width=1.35cm, minimum height=0.58cm,
                 line width=1.2pt},
    pipe/.style={-{Latex[length=2.4mm,width=2mm]}, line width=1pt, draw=nankai!80!black, rounded corners=2pt},
    tag/.style={font=\tiny\bfseries, text=nankai!88!black}
  ]
    % ---- top branch ----
    \node[io]   (img)  at (0,1.15) {\textbf{原图 + bbox}};
    \node[sam]  (sam0) at (2.2,1.15) {\textbf{SAM2}\\[0.5pt]{\tiny bbox-only}};
    \node[cand] (c0)   at (4.2,1.15) {\tiny 3 candidates};

    % ---- bottom branch ----
    \node[io]   (img2) at (0,0) {\textbf{原图 + bbox}\\[0.5pt]{\tiny + positive points}};
    \node[sam]  (sam1) at (2.2,0) {\textbf{SAM2}\\[0.5pt]{\tiny bbox + pos}};
    \node[cand] (c1)   at (4.2,0) {\tiny up to 3 candidates};

    % ---- merge into rerank ----
    \node[ours] (rerank) at (6.3,0.575) {\textbf{Ensemble}\\[0.5pt]{\tiny + Rerank}};
    \node[best] (out)    at (8.3,0.575) {\textbf{最优 mask}\\[0.5pt]{\tiny $\alpha$}};

    % ---- arrows ----
    \draw[pipe] (img)  -- (sam0);
    \draw[pipe] (sam0) -- (c0);
    \draw[pipe] (img2) -- (sam1);
    \draw[pipe] (sam1) -- (c1);
    \draw[pipe] (c0.east) -| (rerank.north);
    \draw[pipe] (c1.east) -| (rerank.south);
    \draw[pipe] (rerank) -- (out);

    % ---- grouping bands ----
    \begin{scope}[on background layer]
      \node[fit=(img)(sam0)(c0)(img2)(sam1)(c1), rounded corners=6pt,
            draw=accentblue!45, line width=0.8pt, fill=accentblue!5, inner sep=8pt] (grpA) {};
      \node[fit=(rerank)(out), rounded corners=6pt,
            draw=nankai!52, line width=0.8pt, fill=nankai!6, inner sep=8pt] (grpB) {};
    \end{scope}
    \node[tag, anchor=south west] at ([xshift=6pt,yshift=-1pt]grpA.north west)
      {多策略候选生成 · Multi-prompt generation};
    \node[tag, anchor=south east] at ([xshift=-6pt,yshift=-1pt]grpB.north east)
      {置信度重排序 · Confidence-aware rerank};
  \end{tikzpicture}}

  \vspace{0.04em}
  \begin{columns}[T,onlytextwidth]
    \column{0.24\textwidth}
      \ccard{nankai}{0.95cm}{%
      {\fontsize{6}{7}\selectfont\textbf{\textcolor{nankai!88!black}{重排序公式}}}\\[1pt]
      {\fontsize{5.5}{6.5}\selectfont $S(m)=0.55\cdot\text{IoU}(m)+0.25\cdot\text{conn}(m)+0.20\cdot\text{bbox\_cov}(m)$}}
    \column{0.25\textwidth}
      \ccard{accentblue}{0.95cm}{%
      {\fontsize{6}{7}\selectfont\textbf{\textcolor{accentblue!85!black}{SAM2 IoU (55\%)}}}\\[1pt]
      {\fontsize{5.5}{6.5}\selectfont 模型自身的 mask 质量置信度；最可靠信号}}
    \column{0.25\textwidth}
      \ccard{softorange}{0.95cm}{%
      {\fontsize{6}{7}\selectfont\textbf{\textcolor{softorange!75!black}{连通性 (25\%)}}}\\[1pt]
      {\fontsize{5.5}{6.5}\selectfont 连通分量数衡量 mask 碎片化；惩罚零散碎片}}
    \column{0.26\textwidth}
      \ccard{softgreen}{0.95cm}{%
      {\fontsize{6}{7}\selectfont\textbf{\textcolor{softgreen!60!black}{bbox 覆盖 (20\%)}}}\\[1pt]
      {\fontsize{5.5}{6.5}\selectfont mask 应填充 bbox 的 10--70\%；偏离的压低}}
  \end{columns}
\end{frame}

% ==================== 第 3 页：多尺度推理 ====================
\begin{frame}[t]{3. 核心优化二：多尺度推理}
  \scriptsize
  \linespread{0.92}\selectfont

  \notestrip{deepred}{\scriptsize
  \textbf{动机：}单尺度 SAM2 推理对图像分辨率敏感，边界处可能产生偏差。我们引入 \hilite{多尺度推理}：在多个尺度分别运行 SAM2，融合结果以降低单尺度偏差。}

  \vspace{0.70em}
  \begin{columns}[T,onlytextwidth]
    \column{0.50\textwidth}
      \tcard{accentblue}{3.20cm}{多尺度推理原理}{%
        \begin{itemize}
          \setlength{\itemsep}{0.08em}
          \item 在 3 个尺度 $\{0.75\times, 1.0\times, 1.25\times\}$ 分别运行 SAM2；
          \item 将各尺度 mask resize 回原图分辨率；
          \item 平均融合：$\alpha_{\text{fused}}=\frac{1}{3}\sum_{s}\alpha_s$；
          \item 单一尺度的边界偏差被多数投票稀释；
          \item 不增加模型参数，纯推理时后处理。
        \end{itemize}
        \vspace{0.10em}
        \textbf{与 Ensemble 叠加：}\hilite{每尺度各自 ensemble rerank}，再跨尺度平均 —— 两阶段互补，进一步提升鲁棒性。}
    \column{0.50\textwidth}
      \tcard{nankai}{3.20cm}{直观理解}{%
        \begin{itemize}
          \setlength{\itemsep}{0.10em}
          \item \textbf{放大（1.25×）：}捕捉更精细的边界细节；
          \item \textbf{原始（1.0×）：}保持模型训练分布下的最优精度；
          \item \textbf{缩小（0.75×）：}获得更大感受野，减少局部歧义。
        \end{itemize}
        \vspace{0.15em}
        三个尺度的 mask 经 resize + 平均后，\hilite{单尺度的随机偏差被平滑}，类似集成学习中 bagging 的思路。

        \vspace{0.12em}
        \textbf{代价：}推理时间约 $\times$3（3 个尺度），但与 ensemble 组合时收益大于开销。}
  \end{columns}
\end{frame}

% ==================== 第 4 页：消融实验设定 ====================
\begin{frame}[t]{4. 消融实验设定}
  \footnotesize
  \linespread{0.92}\selectfont

  \begin{columns}[T,onlytextwidth]
    \column{0.48\textwidth}
      \tcard{softgray}{2.60cm}{数据集与模型}{%
        \begin{itemize}
          \setlength{\itemsep}{0.08em}
          \item \textbf{数据集：}MicroMat-3K valid706 子集
          \begin{itemize}
            \item 706 张图像，64 张基础图像
            \item 每张配有 official bbox + GT alpha
          \end{itemize}
          \item \textbf{模型：}SAM2.1 Hiera-Base+（80.8M 参数）
        \end{itemize}
      }
      \vspace{0.15em}
      \tcard{accentblue}{1.90cm}{评测指标}{%
        \begin{itemize}
          \setlength{\itemsep}{0.06em}
          \item \textbf{SAD}（Sum of Absolute Differences）—— 绝对误差总和
          \item \textbf{MSE} —— 均方误差
          \item \textbf{MAE$\times$1000} —— 平均绝对误差（$\times$1000）
          \item \textbf{Boundary SAD} —— 边界带（5px）SAD，衡量边缘质量
        \end{itemize}
      }
    \column{0.52\textwidth}
      \tcard{nankai}{4.65cm}{消融方法设计（8 种组合）}{%
        \textbf{核心思路：}控制变量法，逐项检验各部分贡献。
        \vspace{0.08em}
        \begin{itemize}
          \setlength{\itemsep}{0.05em}
          \item \textbf{bbox\_binary} — 基线：bbox-only 单次推理 + 硬二值化
          \item \textbf{+ guided} — 单路 + guided filter 边缘平滑
          \item \textbf{+ multiscale} — 单路 + 3 尺度融合
          \item \textbf{+ multiscale\_guided} — 单路 + 多尺度 + guided
          \item \textbf{+ ensemble\_rerank} — 双 prompt 候选 + IoU 重排序
          \item \textbf{+ ensemble\_guided} — ensemble + guided
          \item \textbf{+ ensemble\_multiscale} — ensemble + 多尺度
          \item \textbf{+ ensemble\_multiscale\_guided} — 全部叠加
        \end{itemize}
        \vspace{0.08em}
        从 1 到 8 逐步叠加优化，分析各模块的\hilite{独立贡献与交互效应}。}
  \end{columns}
\end{frame}

% ==================== 第 5 页：消融结果 ====================
\begin{frame}[t]{5. 消融实验结果（valid706）}
  \scriptsize
  \linespread{0.92}\selectfont

  \begin{columns}[T,onlytextwidth]
    \column{0.58\textwidth}
      \begin{block}{完整消融对比表}
        \vspace{-0.3em}
        \centering
        {\fontsize{7}{8}\selectfont
        \setlength{\tabcolsep}{2.2pt}
        \renewcommand{\arraystretch}{1.04}
        \resizebox{0.99\linewidth}{!}{%
        \begin{tabular}{lrrrr}
          \toprule
          \textbf{方法} & \textbf{SAD $\downarrow$} & \textbf{MSE $\downarrow$} & \textbf{MAE$\times$1000 $\downarrow$} & \textbf{Bnd-SAD $\downarrow$} \\
          \midrule
          \textbf{ensemble\_guided (最终)} & \textbf{2774.0} & \textbf{0.000747} & \textbf{0.981} & \textbf{1215.5} \\
          ensemble\_multiscale\_guided & 2915.4 & 0.000742 & 1.033 & 1188.8 \\
          ensemble\_rerank & 2930.9 & 0.000843 & 1.037 & 1380.1 \\
          ensemble\_multiscale & 2997.4 & 0.000796 & 1.062 & 1279.2 \\
          multiscale\_guided & 3145.3 & 0.000828 & 1.089 & 1249.0 \\
          guided & 3161.9 & 0.000845 & 1.095 & 1271.2 \\
          multiscale & 3228.3 & 0.000894 & 1.119 & 1340.7 \\
          \textbf{bbox\_binary (基线)} & \textbf{3313.5} & \textbf{0.000956} & \textbf{1.149} & \textbf{1431.5} \\
          \bottomrule
        \end{tabular}}}
      \end{block}
    \column{0.42\textwidth}
      \tcard{softgreen}{2.05cm}{关键发现}{%
        \begin{itemize}
          \setlength{\itemsep}{0.08em}
          \item \hilite{Ensemble rerank 单点最有效}（SAD$\downarrow$11.5\% vs 基线）
          \item Guided filter 对 Boundary 改善显著（1432$\to$1271）
          \item Multiscale 单独增益有限（3313$\to$3228）
          \item 三者叠加达最优：SAD$\downarrow$16.3\%，MAE$\downarrow$14.6\%
        \end{itemize}
      }

      \vspace{0.25em}
      \tcard{nankai}{1.30cm}{最终提交方法}{%
        \centering
        Ensemble Rerank + Guided Filter}

      \vspace{0.6em}
      \begin{columns}[T,onlytextwidth]
        \column{0.33\textwidth}
          \metricchip{softgreen}{SAD 改善{\tiny\ (vs 基线)}}{$-16.3\%$}
        \column{0.33\textwidth}
          \metricchip{softgreen}{MAE 改善{\tiny\ (vs 基线)}}{$-14.6\%$}
        \column{0.34\textwidth}
          \metricchip{nankai}{Boundary 改善{\tiny\ (vs 基线)}}{$-15.1\%$}
      \end{columns}
  \end{columns}
\end{frame}

% ==================== 第 6 页：可视化案例一 ====================
\begin{frame}[t]{6. 可视化案例一：大型目标（汽车）}
  \footnotesize
  \linespread{0.92}\selectfont

  \vspace{-0.2em}
  \centering
  \fcolorbox{black!25}{white}{%
    \includegraphics[width=0.92\linewidth,height=0.60\textheight,keepaspectratio]{pic/sam2_case1_compare.png}}

  \vspace{0.12em}
  \begin{columns}[T,onlytextwidth]
    \column{0.50\textwidth}
      \tcard{deepred}{1.10cm}{Baseline 问题}{%
        \small SAD=13559。车轮和车窗边缘有明显的 mask 溢出，导致前景区域偏大。}
    \column{0.50\textwidth}
      \tcard{softgreen}{1.10cm}{优化后改善}{%
        \small SAD=9010（$\downarrow$33.5\%）。边界更贴合车身轮廓，车窗和轮毂区域 mask 精度提升。}
  \end{columns}
\end{frame}

% ==================== 第 7 页：可视化案例二 ====================
\begin{frame}[t]{7. 可视化案例二：小型目标}
  \footnotesize
  \linespread{0.92}\selectfont

  \vspace{-0.2em}
  \centering
  \fcolorbox{black!25}{white}{%
    \includegraphics[width=0.80\linewidth,height=0.60\textheight,keepaspectratio]{pic/sam2_case2_compare.png}}

  \vspace{0.12em}
  \begin{columns}[T,onlytextwidth]
    \column{0.50\textwidth}
      \tcard{deepred}{1.10cm}{Baseline 问题}{%
        \small SAD=3843。微小目标（bbox 仅 117$\times$133）容易误选背景区域，mask 几乎全错。}
    \column{0.50\textwidth}
      \tcard{softgreen}{1.10cm}{优化后改善}{%
        \small SAD=348（$\downarrow$90.9\%）。Ensemble 重排序有效纠正了候选选择错误，精确捕获目标。}
  \end{columns}
\end{frame}

% ==================== 第 8 页：可视化案例三 ====================
\begin{frame}[t]{8. 可视化案例三：细结构物体}
  \footnotesize
  \linespread{0.92}\selectfont

  \vspace{-0.2em}
  \centering
  \fcolorbox{black!25}{white}{%
    \includegraphics[width=0.80\linewidth,height=0.60\textheight,keepaspectratio]{pic/sam2_case3_compare.png}}

  \vspace{0.12em}
  \begin{columns}[T,onlytextwidth]
    \column{0.50\textwidth}
      \tcard{deepred}{1.10cm}{Baseline 问题}{%
        \small SAD=2562。细结构处有 mask 碎片和空洞，边界不完整。}
    \column{0.50\textwidth}
      \tcard{softgreen}{1.10cm}{优化后改善}{%
        \small SAD=374（$\downarrow$85.4\%）。Mask 更完整，碎片显著减少，细边缘更清晰。}
  \end{columns}
\end{frame}

% ---- 本节结束 ----
 引入，
%    继承主文件导言区的全部宏 / 配色 / 字体，无需自带 \documentclass 或 \begin{document}。
%  · 内容：模型选型 → Prompt Ensemble + Rerank → 多尺度融合 + 消融 → 可视化
% ============================================================================
\section{上游分割}

% ==================== 第 1 页：模型选型与系统定位 ====================
\begin{frame}[t]{1. 上游模型选型：为什么用 SAM2？}
  \footnotesize
  \linespread{0.90}\selectfont

  \begin{columns}[T,onlytextwidth]
    \column{0.50\textwidth}
      \tcard{softorange}{2.00cm}{定位需求分析}{%
        \textbf{系统需求：}给定图像，输出 bbox / points，供 ZIM 做精细 matting。\\
        \textbf{关键约束：}
        \begin{itemize}
          \setlength{\itemsep}{0.02em}
          \item \hilite{zero-training} — 不微调模型参数；
          \item 推理高效 — 单卡秒级；
          \item bbox prompt 原生支持 — 与下游接口统一。
        \end{itemize}
        \vspace{0.02em}
        \textbf{不选 SAM3：}模型更大、推理慢约 2--3 倍，bbox-only 提升有限。\\
        \textbf{不选 Grounded-SAM-2：}额外引入 text encoder，增加显存开销。
      }
    \column{0.50\textwidth}
      \tcard{accentblue}{2.00cm}{SAM2 适配性}{%
        \begin{itemize}
          \setlength{\itemsep}{0.04em}
          \item \textbf{Hiera 架构} — 轻量高效，Hiera-Base+ 仅 80.8M 参数；
          \item \textbf{native bbox + points} — 无需适配层；
          \item \textbf{多候选 mask + IoU 评分} — 为 rerank 提供置信度；
          \item \textbf{实测性能} — A10 上约 1.16 it/s（bbox-only）；
          \item 社区活跃，checkpoint 可下载。
        \end{itemize}
      }
  \end{columns}

  \vspace{0.15em}
  \notestrip{accentblue}{\centering\scriptsize
  \textbf{定位：}SAM2 负责回答 \hilite{``目标在哪''} — 输出 bbox + points + 候选 mask；精细软 alpha 留给下游 ZIM + TTA-GF + BSP。}
\end{frame}

% ==================== 第 2 页：Prompt Ensemble + Mask Rerank ====================
\begin{frame}[t]{2. 核心优化一：Prompt Ensemble 与 IoU 引导重排序}
  \scriptsize
  \linespread{0.90}\selectfont

  \notestrip{deepred}{\tiny
  \textbf{动机：}SAM2 单次 bbox 推理易受局部纹理干扰，可能选到次优 mask。我们提出 \hilite{Prompt Ensemble + IoU 引导重排序}：多组 prompt 策略生成候选 mask，再利用 SAM2 内置 IoU 与结构先验选出最优。}

  \vspace{0.35em}
  \centering
  \resizebox{0.92\linewidth}{!}{%
  \begin{tikzpicture}[
    x=1cm, y=1cm,
    >={Latex[length=2.5mm,width=2mm]},
    every node/.style={font=\tiny},
    softsh/.style={blur shadow={shadow blur steps=5, shadow xshift=0.4pt,
                   shadow yshift=-0.5pt, shadow opacity=24, shadow blur radius=1.5pt}},
    nodebase/.style={softsh, rounded corners=3pt, line width=0.9pt, align=center, inner sep=3pt},
    io/.style={nodebase, draw=black!50, fill=black!5, minimum width=1.5cm, minimum height=0.58cm},
    sam/.style={nodebase, draw=accentblue!82!black, fill=accentblue!14, minimum width=1.55cm, minimum height=0.58cm},
    cand/.style={nodebase, draw=softorange!72!black, fill=softorange!12, minimum width=1.2cm, minimum height=0.58cm},
    ours/.style={nodebase, draw=nankai!88!black, fill=nankai!14, minimum width=1.35cm, minimum height=0.58cm},
    best/.style={nodebase, draw=softgreen!72!black, fill=softgreen!16, minimum width=1.35cm, minimum height=0.58cm,
                 line width=1.2pt},
    pipe/.style={-{Latex[length=2.4mm,width=2mm]}, line width=1pt, draw=nankai!80!black, rounded corners=2pt},
    tag/.style={font=\tiny\bfseries, text=nankai!88!black}
  ]
    % ---- top branch ----
    \node[io]   (img)  at (0,1.15) {\textbf{原图 + bbox}};
    \node[sam]  (sam0) at (2.2,1.15) {\textbf{SAM2}\\[0.5pt]{\tiny bbox-only}};
    \node[cand] (c0)   at (4.2,1.15) {\tiny 3 candidates};

    % ---- bottom branch ----
    \node[io]   (img2) at (0,0) {\textbf{原图 + bbox}\\[0.5pt]{\tiny + positive points}};
    \node[sam]  (sam1) at (2.2,0) {\textbf{SAM2}\\[0.5pt]{\tiny bbox + pos}};
    \node[cand] (c1)   at (4.2,0) {\tiny up to 3 candidates};

    % ---- merge into rerank ----
    \node[ours] (rerank) at (6.3,0.575) {\textbf{Ensemble}\\[0.5pt]{\tiny + Rerank}};
    \node[best] (out)    at (8.3,0.575) {\textbf{最优 mask}\\[0.5pt]{\tiny $\alpha$}};

    % ---- arrows ----
    \draw[pipe] (img)  -- (sam0);
    \draw[pipe] (sam0) -- (c0);
    \draw[pipe] (img2) -- (sam1);
    \draw[pipe] (sam1) -- (c1);
    \draw[pipe] (c0.east) -| (rerank.north);
    \draw[pipe] (c1.east) -| (rerank.south);
    \draw[pipe] (rerank) -- (out);

    % ---- grouping bands ----
    \begin{scope}[on background layer]
      \node[fit=(img)(sam0)(c0)(img2)(sam1)(c1), rounded corners=6pt,
            draw=accentblue!45, line width=0.8pt, fill=accentblue!5, inner sep=8pt] (grpA) {};
      \node[fit=(rerank)(out), rounded corners=6pt,
            draw=nankai!52, line width=0.8pt, fill=nankai!6, inner sep=8pt] (grpB) {};
    \end{scope}
    \node[tag, anchor=south west] at ([xshift=6pt,yshift=-1pt]grpA.north west)
      {多策略候选生成 · Multi-prompt generation};
    \node[tag, anchor=south east] at ([xshift=-6pt,yshift=-1pt]grpB.north east)
      {置信度重排序 · Confidence-aware rerank};
  \end{tikzpicture}}

  \vspace{0.04em}
  \begin{columns}[T,onlytextwidth]
    \column{0.24\textwidth}
      \ccard{nankai}{0.95cm}{%
      {\fontsize{6}{7}\selectfont\textbf{\textcolor{nankai!88!black}{重排序公式}}}\\[1pt]
      {\fontsize{5.5}{6.5}\selectfont $S(m)=0.55\cdot\text{IoU}(m)+0.25\cdot\text{conn}(m)+0.20\cdot\text{bbox\_cov}(m)$}}
    \column{0.25\textwidth}
      \ccard{accentblue}{0.95cm}{%
      {\fontsize{6}{7}\selectfont\textbf{\textcolor{accentblue!85!black}{SAM2 IoU (55\%)}}}\\[1pt]
      {\fontsize{5.5}{6.5}\selectfont 模型自身的 mask 质量置信度；最可靠信号}}
    \column{0.25\textwidth}
      \ccard{softorange}{0.95cm}{%
      {\fontsize{6}{7}\selectfont\textbf{\textcolor{softorange!75!black}{连通性 (25\%)}}}\\[1pt]
      {\fontsize{5.5}{6.5}\selectfont 连通分量数衡量 mask 碎片化；惩罚零散碎片}}
    \column{0.26\textwidth}
      \ccard{softgreen}{0.95cm}{%
      {\fontsize{6}{7}\selectfont\textbf{\textcolor{softgreen!60!black}{bbox 覆盖 (20\%)}}}\\[1pt]
      {\fontsize{5.5}{6.5}\selectfont mask 应填充 bbox 的 10--70\%；偏离的压低}}
  \end{columns}
\end{frame}

% ==================== 第 3 页：多尺度推理 ====================
\begin{frame}[t]{3. 核心优化二：多尺度推理}
  \scriptsize
  \linespread{0.92}\selectfont

  \notestrip{deepred}{\scriptsize
  \textbf{动机：}单尺度 SAM2 推理对图像分辨率敏感，边界处可能产生偏差。我们引入 \hilite{多尺度推理}：在多个尺度分别运行 SAM2，融合结果以降低单尺度偏差。}

  \vspace{0.70em}
  \begin{columns}[T,onlytextwidth]
    \column{0.50\textwidth}
      \tcard{accentblue}{3.20cm}{多尺度推理原理}{%
        \begin{itemize}
          \setlength{\itemsep}{0.08em}
          \item 在 3 个尺度 $\{0.75\times, 1.0\times, 1.25\times\}$ 分别运行 SAM2；
          \item 将各尺度 mask resize 回原图分辨率；
          \item 平均融合：$\alpha_{\text{fused}}=\frac{1}{3}\sum_{s}\alpha_s$；
          \item 单一尺度的边界偏差被多数投票稀释；
          \item 不增加模型参数，纯推理时后处理。
        \end{itemize}
        \vspace{0.10em}
        \textbf{与 Ensemble 叠加：}\hilite{每尺度各自 ensemble rerank}，再跨尺度平均 —— 两阶段互补，进一步提升鲁棒性。}
    \column{0.50\textwidth}
      \tcard{nankai}{3.20cm}{直观理解}{%
        \begin{itemize}
          \setlength{\itemsep}{0.10em}
          \item \textbf{放大（1.25×）：}捕捉更精细的边界细节；
          \item \textbf{原始（1.0×）：}保持模型训练分布下的最优精度；
          \item \textbf{缩小（0.75×）：}获得更大感受野，减少局部歧义。
        \end{itemize}
        \vspace{0.15em}
        三个尺度的 mask 经 resize + 平均后，\hilite{单尺度的随机偏差被平滑}，类似集成学习中 bagging 的思路。

        \vspace{0.12em}
        \textbf{代价：}推理时间约 $\times$3（3 个尺度），但与 ensemble 组合时收益大于开销。}
  \end{columns}
\end{frame}

% ==================== 第 4 页：消融实验设定 ====================
\begin{frame}[t]{4. 消融实验设定}
  \footnotesize
  \linespread{0.92}\selectfont

  \begin{columns}[T,onlytextwidth]
    \column{0.48\textwidth}
      \tcard{softgray}{2.60cm}{数据集与模型}{%
        \begin{itemize}
          \setlength{\itemsep}{0.08em}
          \item \textbf{数据集：}MicroMat-3K valid706 子集
          \begin{itemize}
            \item 706 张图像，64 张基础图像
            \item 每张配有 official bbox + GT alpha
          \end{itemize}
          \item \textbf{模型：}SAM2.1 Hiera-Base+（80.8M 参数）
        \end{itemize}
      }
      \vspace{0.15em}
      \tcard{accentblue}{1.90cm}{评测指标}{%
        \begin{itemize}
          \setlength{\itemsep}{0.06em}
          \item \textbf{SAD}（Sum of Absolute Differences）—— 绝对误差总和
          \item \textbf{MSE} —— 均方误差
          \item \textbf{MAE$\times$1000} —— 平均绝对误差（$\times$1000）
          \item \textbf{Boundary SAD} —— 边界带（5px）SAD，衡量边缘质量
        \end{itemize}
      }
    \column{0.52\textwidth}
      \tcard{nankai}{4.65cm}{消融方法设计（8 种组合）}{%
        \textbf{核心思路：}控制变量法，逐项检验各部分贡献。
        \vspace{0.08em}
        \begin{itemize}
          \setlength{\itemsep}{0.05em}
          \item \textbf{bbox\_binary} — 基线：bbox-only 单次推理 + 硬二值化
          \item \textbf{+ guided} — 单路 + guided filter 边缘平滑
          \item \textbf{+ multiscale} — 单路 + 3 尺度融合
          \item \textbf{+ multiscale\_guided} — 单路 + 多尺度 + guided
          \item \textbf{+ ensemble\_rerank} — 双 prompt 候选 + IoU 重排序
          \item \textbf{+ ensemble\_guided} — ensemble + guided
          \item \textbf{+ ensemble\_multiscale} — ensemble + 多尺度
          \item \textbf{+ ensemble\_multiscale\_guided} — 全部叠加
        \end{itemize}
        \vspace{0.08em}
        从 1 到 8 逐步叠加优化，分析各模块的\hilite{独立贡献与交互效应}。}
  \end{columns}
\end{frame}

% ==================== 第 5 页：消融结果 ====================
\begin{frame}[t]{5. 消融实验结果（valid706）}
  \scriptsize
  \linespread{0.92}\selectfont

  \begin{columns}[T,onlytextwidth]
    \column{0.58\textwidth}
      \begin{block}{完整消融对比表}
        \vspace{-0.3em}
        \centering
        {\fontsize{7}{8}\selectfont
        \setlength{\tabcolsep}{2.2pt}
        \renewcommand{\arraystretch}{1.04}
        \resizebox{0.99\linewidth}{!}{%
        \begin{tabular}{lrrrr}
          \toprule
          \textbf{方法} & \textbf{SAD $\downarrow$} & \textbf{MSE $\downarrow$} & \textbf{MAE$\times$1000 $\downarrow$} & \textbf{Bnd-SAD $\downarrow$} \\
          \midrule
          \textbf{ensemble\_guided (最终)} & \textbf{2774.0} & \textbf{0.000747} & \textbf{0.981} & \textbf{1215.5} \\
          ensemble\_multiscale\_guided & 2915.4 & 0.000742 & 1.033 & 1188.8 \\
          ensemble\_rerank & 2930.9 & 0.000843 & 1.037 & 1380.1 \\
          ensemble\_multiscale & 2997.4 & 0.000796 & 1.062 & 1279.2 \\
          multiscale\_guided & 3145.3 & 0.000828 & 1.089 & 1249.0 \\
          guided & 3161.9 & 0.000845 & 1.095 & 1271.2 \\
          multiscale & 3228.3 & 0.000894 & 1.119 & 1340.7 \\
          \textbf{bbox\_binary (基线)} & \textbf{3313.5} & \textbf{0.000956} & \textbf{1.149} & \textbf{1431.5} \\
          \bottomrule
        \end{tabular}}}
      \end{block}
    \column{0.42\textwidth}
      \tcard{softgreen}{2.05cm}{关键发现}{%
        \begin{itemize}
          \setlength{\itemsep}{0.08em}
          \item \hilite{Ensemble rerank 单点最有效}（SAD$\downarrow$11.5\% vs 基线）
          \item Guided filter 对 Boundary 改善显著（1432$\to$1271）
          \item Multiscale 单独增益有限（3313$\to$3228）
          \item 三者叠加达最优：SAD$\downarrow$16.3\%，MAE$\downarrow$14.6\%
        \end{itemize}
      }

      \vspace{0.25em}
      \tcard{nankai}{1.30cm}{最终提交方法}{%
        \centering
        Ensemble Rerank + Guided Filter}

      \vspace{0.6em}
      \begin{columns}[T,onlytextwidth]
        \column{0.33\textwidth}
          \metricchip{softgreen}{SAD 改善{\tiny\ (vs 基线)}}{$-16.3\%$}
        \column{0.33\textwidth}
          \metricchip{softgreen}{MAE 改善{\tiny\ (vs 基线)}}{$-14.6\%$}
        \column{0.34\textwidth}
          \metricchip{nankai}{Boundary 改善{\tiny\ (vs 基线)}}{$-15.1\%$}
      \end{columns}
  \end{columns}
\end{frame}

% ==================== 第 6 页：可视化案例一 ====================
\begin{frame}[t]{6. 可视化案例一：大型目标（汽车）}
  \footnotesize
  \linespread{0.92}\selectfont

  \vspace{-0.2em}
  \centering
  \fcolorbox{black!25}{white}{%
    \includegraphics[width=0.92\linewidth,height=0.60\textheight,keepaspectratio]{pic/sam2_case1_compare.png}}

  \vspace{0.12em}
  \begin{columns}[T,onlytextwidth]
    \column{0.50\textwidth}
      \tcard{deepred}{1.10cm}{Baseline 问题}{%
        \small SAD=13559。车轮和车窗边缘有明显的 mask 溢出，导致前景区域偏大。}
    \column{0.50\textwidth}
      \tcard{softgreen}{1.10cm}{优化后改善}{%
        \small SAD=9010（$\downarrow$33.5\%）。边界更贴合车身轮廓，车窗和轮毂区域 mask 精度提升。}
  \end{columns}
\end{frame}

% ==================== 第 7 页：可视化案例二 ====================
\begin{frame}[t]{7. 可视化案例二：小型目标}
  \footnotesize
  \linespread{0.92}\selectfont

  \vspace{-0.2em}
  \centering
  \fcolorbox{black!25}{white}{%
    \includegraphics[width=0.80\linewidth,height=0.60\textheight,keepaspectratio]{pic/sam2_case2_compare.png}}

  \vspace{0.12em}
  \begin{columns}[T,onlytextwidth]
    \column{0.50\textwidth}
      \tcard{deepred}{1.10cm}{Baseline 问题}{%
        \small SAD=3843。微小目标（bbox 仅 117$\times$133）容易误选背景区域，mask 几乎全错。}
    \column{0.50\textwidth}
      \tcard{softgreen}{1.10cm}{优化后改善}{%
        \small SAD=348（$\downarrow$90.9\%）。Ensemble 重排序有效纠正了候选选择错误，精确捕获目标。}
  \end{columns}
\end{frame}

% ==================== 第 8 页：可视化案例三 ====================
\begin{frame}[t]{8. 可视化案例三：细结构物体}
  \footnotesize
  \linespread{0.92}\selectfont

  \vspace{-0.2em}
  \centering
  \fcolorbox{black!25}{white}{%
    \includegraphics[width=0.80\linewidth,height=0.60\textheight,keepaspectratio]{pic/sam2_case3_compare.png}}

  \vspace{0.12em}
  \begin{columns}[T,onlytextwidth]
    \column{0.50\textwidth}
      \tcard{deepred}{1.10cm}{Baseline 问题}{%
        \small SAD=2562。细结构处有 mask 碎片和空洞，边界不完整。}
    \column{0.50\textwidth}
      \tcard{softgreen}{1.10cm}{优化后改善}{%
        \small SAD=374（$\downarrow$85.4\%）。Mask 更完整，碎片显著减少，细边缘更清晰。}
  \end{columns}
\end{frame}

% ---- 本节结束 ----
 引入，
%    继承主文件导言区的全部宏 / 配色 / 字体，无需自带 \documentclass 或 \begin{document}。
%  · 内容：模型选型 → Prompt Ensemble + Rerank → 多尺度融合 + 消融 → 可视化
% ============================================================================
\section{上游分割}

% ==================== 第 1 页：模型选型与系统定位 ====================
\begin{frame}[t]{1. 上游模型选型：为什么用 SAM2？}
  \footnotesize
  \linespread{0.90}\selectfont

  \begin{columns}[T,onlytextwidth]
    \column{0.50\textwidth}
      \tcard{softorange}{2.00cm}{定位需求分析}{%
        \textbf{系统需求：}给定图像，输出 bbox / points，供 ZIM 做精细 matting。\\
        \textbf{关键约束：}
        \begin{itemize}
          \setlength{\itemsep}{0.02em}
          \item \hilite{zero-training} — 不微调模型参数；
          \item 推理高效 — 单卡秒级；
          \item bbox prompt 原生支持 — 与下游接口统一。
        \end{itemize}
        \vspace{0.02em}
        \textbf{不选 SAM3：}模型更大、推理慢约 2--3 倍，bbox-only 提升有限。\\
        \textbf{不选 Grounded-SAM-2：}额外引入 text encoder，增加显存开销。
      }
    \column{0.50\textwidth}
      \tcard{accentblue}{2.00cm}{SAM2 适配性}{%
        \begin{itemize}
          \setlength{\itemsep}{0.04em}
          \item \textbf{Hiera 架构} — 轻量高效，Hiera-Base+ 仅 80.8M 参数；
          \item \textbf{native bbox + points} — 无需适配层；
          \item \textbf{多候选 mask + IoU 评分} — 为 rerank 提供置信度；
          \item \textbf{实测性能} — A10 上约 1.16 it/s（bbox-only）；
          \item 社区活跃，checkpoint 可下载。
        \end{itemize}
      }
  \end{columns}

  \vspace{0.15em}
  \notestrip{accentblue}{\centering\scriptsize
  \textbf{定位：}SAM2 负责回答 \hilite{``目标在哪''} — 输出 bbox + points + 候选 mask；精细软 alpha 留给下游 ZIM + TTA-GF + BSP。}
\end{frame}

% ==================== 第 2 页：Prompt Ensemble + Mask Rerank ====================
\begin{frame}[t]{2. 核心优化一：Prompt Ensemble 与 IoU 引导重排序}
  \scriptsize
  \linespread{0.90}\selectfont

  \notestrip{deepred}{\tiny
  \textbf{动机：}SAM2 单次 bbox 推理易受局部纹理干扰，可能选到次优 mask。我们提出 \hilite{Prompt Ensemble + IoU 引导重排序}：多组 prompt 策略生成候选 mask，再利用 SAM2 内置 IoU 与结构先验选出最优。}

  \vspace{0.35em}
  \centering
  \resizebox{0.92\linewidth}{!}{%
  \begin{tikzpicture}[
    x=1cm, y=1cm,
    >={Latex[length=2.5mm,width=2mm]},
    every node/.style={font=\tiny},
    softsh/.style={blur shadow={shadow blur steps=5, shadow xshift=0.4pt,
                   shadow yshift=-0.5pt, shadow opacity=24, shadow blur radius=1.5pt}},
    nodebase/.style={softsh, rounded corners=3pt, line width=0.9pt, align=center, inner sep=3pt},
    io/.style={nodebase, draw=black!50, fill=black!5, minimum width=1.5cm, minimum height=0.58cm},
    sam/.style={nodebase, draw=accentblue!82!black, fill=accentblue!14, minimum width=1.55cm, minimum height=0.58cm},
    cand/.style={nodebase, draw=softorange!72!black, fill=softorange!12, minimum width=1.2cm, minimum height=0.58cm},
    ours/.style={nodebase, draw=nankai!88!black, fill=nankai!14, minimum width=1.35cm, minimum height=0.58cm},
    best/.style={nodebase, draw=softgreen!72!black, fill=softgreen!16, minimum width=1.35cm, minimum height=0.58cm,
                 line width=1.2pt},
    pipe/.style={-{Latex[length=2.4mm,width=2mm]}, line width=1pt, draw=nankai!80!black, rounded corners=2pt},
    tag/.style={font=\tiny\bfseries, text=nankai!88!black}
  ]
    % ---- top branch ----
    \node[io]   (img)  at (0,1.15) {\textbf{原图 + bbox}};
    \node[sam]  (sam0) at (2.2,1.15) {\textbf{SAM2}\\[0.5pt]{\tiny bbox-only}};
    \node[cand] (c0)   at (4.2,1.15) {\tiny 3 candidates};

    % ---- bottom branch ----
    \node[io]   (img2) at (0,0) {\textbf{原图 + bbox}\\[0.5pt]{\tiny + positive points}};
    \node[sam]  (sam1) at (2.2,0) {\textbf{SAM2}\\[0.5pt]{\tiny bbox + pos}};
    \node[cand] (c1)   at (4.2,0) {\tiny up to 3 candidates};

    % ---- merge into rerank ----
    \node[ours] (rerank) at (6.3,0.575) {\textbf{Ensemble}\\[0.5pt]{\tiny + Rerank}};
    \node[best] (out)    at (8.3,0.575) {\textbf{最优 mask}\\[0.5pt]{\tiny $\alpha$}};

    % ---- arrows ----
    \draw[pipe] (img)  -- (sam0);
    \draw[pipe] (sam0) -- (c0);
    \draw[pipe] (img2) -- (sam1);
    \draw[pipe] (sam1) -- (c1);
    \draw[pipe] (c0.east) -| (rerank.north);
    \draw[pipe] (c1.east) -| (rerank.south);
    \draw[pipe] (rerank) -- (out);

    % ---- grouping bands ----
    \begin{scope}[on background layer]
      \node[fit=(img)(sam0)(c0)(img2)(sam1)(c1), rounded corners=6pt,
            draw=accentblue!45, line width=0.8pt, fill=accentblue!5, inner sep=8pt] (grpA) {};
      \node[fit=(rerank)(out), rounded corners=6pt,
            draw=nankai!52, line width=0.8pt, fill=nankai!6, inner sep=8pt] (grpB) {};
    \end{scope}
    \node[tag, anchor=south west] at ([xshift=6pt,yshift=-1pt]grpA.north west)
      {多策略候选生成 · Multi-prompt generation};
    \node[tag, anchor=south east] at ([xshift=-6pt,yshift=-1pt]grpB.north east)
      {置信度重排序 · Confidence-aware rerank};
  \end{tikzpicture}}

  \vspace{0.04em}
  \begin{columns}[T,onlytextwidth]
    \column{0.24\textwidth}
      \ccard{nankai}{0.95cm}{%
      {\fontsize{6}{7}\selectfont\textbf{\textcolor{nankai!88!black}{重排序公式}}}\\[1pt]
      {\fontsize{5.5}{6.5}\selectfont $S(m)=0.55\cdot\text{IoU}(m)+0.25\cdot\text{conn}(m)+0.20\cdot\text{bbox\_cov}(m)$}}
    \column{0.25\textwidth}
      \ccard{accentblue}{0.95cm}{%
      {\fontsize{6}{7}\selectfont\textbf{\textcolor{accentblue!85!black}{SAM2 IoU (55\%)}}}\\[1pt]
      {\fontsize{5.5}{6.5}\selectfont 模型自身的 mask 质量置信度；最可靠信号}}
    \column{0.25\textwidth}
      \ccard{softorange}{0.95cm}{%
      {\fontsize{6}{7}\selectfont\textbf{\textcolor{softorange!75!black}{连通性 (25\%)}}}\\[1pt]
      {\fontsize{5.5}{6.5}\selectfont 连通分量数衡量 mask 碎片化；惩罚零散碎片}}
    \column{0.26\textwidth}
      \ccard{softgreen}{0.95cm}{%
      {\fontsize{6}{7}\selectfont\textbf{\textcolor{softgreen!60!black}{bbox 覆盖 (20\%)}}}\\[1pt]
      {\fontsize{5.5}{6.5}\selectfont mask 应填充 bbox 的 10--70\%；偏离的压低}}
  \end{columns}
\end{frame}

% ==================== 第 3 页：多尺度推理 ====================
\begin{frame}[t]{3. 核心优化二：多尺度推理}
  \scriptsize
  \linespread{0.92}\selectfont

  \notestrip{deepred}{\scriptsize
  \textbf{动机：}单尺度 SAM2 推理对图像分辨率敏感，边界处可能产生偏差。我们引入 \hilite{多尺度推理}：在多个尺度分别运行 SAM2，融合结果以降低单尺度偏差。}

  \vspace{0.70em}
  \begin{columns}[T,onlytextwidth]
    \column{0.50\textwidth}
      \tcard{accentblue}{3.20cm}{多尺度推理原理}{%
        \begin{itemize}
          \setlength{\itemsep}{0.08em}
          \item 在 3 个尺度 $\{0.75\times, 1.0\times, 1.25\times\}$ 分别运行 SAM2；
          \item 将各尺度 mask resize 回原图分辨率；
          \item 平均融合：$\alpha_{\text{fused}}=\frac{1}{3}\sum_{s}\alpha_s$；
          \item 单一尺度的边界偏差被多数投票稀释；
          \item 不增加模型参数，纯推理时后处理。
        \end{itemize}
        \vspace{0.10em}
        \textbf{与 Ensemble 叠加：}\hilite{每尺度各自 ensemble rerank}，再跨尺度平均 —— 两阶段互补，进一步提升鲁棒性。}
    \column{0.50\textwidth}
      \tcard{nankai}{3.20cm}{直观理解}{%
        \begin{itemize}
          \setlength{\itemsep}{0.10em}
          \item \textbf{放大（1.25×）：}捕捉更精细的边界细节；
          \item \textbf{原始（1.0×）：}保持模型训练分布下的最优精度；
          \item \textbf{缩小（0.75×）：}获得更大感受野，减少局部歧义。
        \end{itemize}
        \vspace{0.15em}
        三个尺度的 mask 经 resize + 平均后，\hilite{单尺度的随机偏差被平滑}，类似集成学习中 bagging 的思路。

        \vspace{0.12em}
        \textbf{代价：}推理时间约 $\times$3（3 个尺度），但与 ensemble 组合时收益大于开销。}
  \end{columns}
\end{frame}

% ==================== 第 4 页：消融实验设定 ====================
\begin{frame}[t]{4. 消融实验设定}
  \footnotesize
  \linespread{0.92}\selectfont

  \begin{columns}[T,onlytextwidth]
    \column{0.48\textwidth}
      \tcard{softgray}{2.60cm}{数据集与模型}{%
        \begin{itemize}
          \setlength{\itemsep}{0.08em}
          \item \textbf{数据集：}MicroMat-3K valid706 子集
          \begin{itemize}
            \item 706 张图像，64 张基础图像
            \item 每张配有 official bbox + GT alpha
          \end{itemize}
          \item \textbf{模型：}SAM2.1 Hiera-Base+（80.8M 参数）
        \end{itemize}
      }
      \vspace{0.15em}
      \tcard{accentblue}{1.90cm}{评测指标}{%
        \begin{itemize}
          \setlength{\itemsep}{0.06em}
          \item \textbf{SAD}（Sum of Absolute Differences）—— 绝对误差总和
          \item \textbf{MSE} —— 均方误差
          \item \textbf{MAE$\times$1000} —— 平均绝对误差（$\times$1000）
          \item \textbf{Boundary SAD} —— 边界带（5px）SAD，衡量边缘质量
        \end{itemize}
      }
    \column{0.52\textwidth}
      \tcard{nankai}{4.65cm}{消融方法设计（8 种组合）}{%
        \textbf{核心思路：}控制变量法，逐项检验各部分贡献。
        \vspace{0.08em}
        \begin{itemize}
          \setlength{\itemsep}{0.05em}
          \item \textbf{bbox\_binary} — 基线：bbox-only 单次推理 + 硬二值化
          \item \textbf{+ guided} — 单路 + guided filter 边缘平滑
          \item \textbf{+ multiscale} — 单路 + 3 尺度融合
          \item \textbf{+ multiscale\_guided} — 单路 + 多尺度 + guided
          \item \textbf{+ ensemble\_rerank} — 双 prompt 候选 + IoU 重排序
          \item \textbf{+ ensemble\_guided} — ensemble + guided
          \item \textbf{+ ensemble\_multiscale} — ensemble + 多尺度
          \item \textbf{+ ensemble\_multiscale\_guided} — 全部叠加
        \end{itemize}
        \vspace{0.08em}
        从 1 到 8 逐步叠加优化，分析各模块的\hilite{独立贡献与交互效应}。}
  \end{columns}
\end{frame}

% ==================== 第 5 页：消融结果 ====================
\begin{frame}[t]{5. 消融实验结果（valid706）}
  \scriptsize
  \linespread{0.92}\selectfont

  \begin{columns}[T,onlytextwidth]
    \column{0.58\textwidth}
      \begin{block}{完整消融对比表}
        \vspace{-0.3em}
        \centering
        {\fontsize{7}{8}\selectfont
        \setlength{\tabcolsep}{2.2pt}
        \renewcommand{\arraystretch}{1.04}
        \resizebox{0.99\linewidth}{!}{%
        \begin{tabular}{lrrrr}
          \toprule
          \textbf{方法} & \textbf{SAD $\downarrow$} & \textbf{MSE $\downarrow$} & \textbf{MAE$\times$1000 $\downarrow$} & \textbf{Bnd-SAD $\downarrow$} \\
          \midrule
          \textbf{ensemble\_guided (最终)} & \textbf{2774.0} & \textbf{0.000747} & \textbf{0.981} & \textbf{1215.5} \\
          ensemble\_multiscale\_guided & 2915.4 & 0.000742 & 1.033 & 1188.8 \\
          ensemble\_rerank & 2930.9 & 0.000843 & 1.037 & 1380.1 \\
          ensemble\_multiscale & 2997.4 & 0.000796 & 1.062 & 1279.2 \\
          multiscale\_guided & 3145.3 & 0.000828 & 1.089 & 1249.0 \\
          guided & 3161.9 & 0.000845 & 1.095 & 1271.2 \\
          multiscale & 3228.3 & 0.000894 & 1.119 & 1340.7 \\
          \textbf{bbox\_binary (基线)} & \textbf{3313.5} & \textbf{0.000956} & \textbf{1.149} & \textbf{1431.5} \\
          \bottomrule
        \end{tabular}}}
      \end{block}
    \column{0.42\textwidth}
      \tcard{softgreen}{2.05cm}{关键发现}{%
        \begin{itemize}
          \setlength{\itemsep}{0.08em}
          \item \hilite{Ensemble rerank 单点最有效}（SAD$\downarrow$11.5\% vs 基线）
          \item Guided filter 对 Boundary 改善显著（1432$\to$1271）
          \item Multiscale 单独增益有限（3313$\to$3228）
          \item 三者叠加达最优：SAD$\downarrow$16.3\%，MAE$\downarrow$14.6\%
        \end{itemize}
      }

      \vspace{0.25em}
      \tcard{nankai}{1.30cm}{最终提交方法}{%
        \centering
        Ensemble Rerank + Guided Filter}

      \vspace{0.6em}
      \begin{columns}[T,onlytextwidth]
        \column{0.33\textwidth}
          \metricchip{softgreen}{SAD 改善{\tiny\ (vs 基线)}}{$-16.3\%$}
        \column{0.33\textwidth}
          \metricchip{softgreen}{MAE 改善{\tiny\ (vs 基线)}}{$-14.6\%$}
        \column{0.34\textwidth}
          \metricchip{nankai}{Boundary 改善{\tiny\ (vs 基线)}}{$-15.1\%$}
      \end{columns}
  \end{columns}
\end{frame}

% ==================== 第 6 页：可视化案例一 ====================
\begin{frame}[t]{6. 可视化案例一：大型目标（汽车）}
  \footnotesize
  \linespread{0.92}\selectfont

  \vspace{-0.2em}
  \centering
  \fcolorbox{black!25}{white}{%
    \includegraphics[width=0.92\linewidth,height=0.60\textheight,keepaspectratio]{pic/sam2_case1_compare.png}}

  \vspace{0.12em}
  \begin{columns}[T,onlytextwidth]
    \column{0.50\textwidth}
      \tcard{deepred}{1.10cm}{Baseline 问题}{%
        \small SAD=13559。车轮和车窗边缘有明显的 mask 溢出，导致前景区域偏大。}
    \column{0.50\textwidth}
      \tcard{softgreen}{1.10cm}{优化后改善}{%
        \small SAD=9010（$\downarrow$33.5\%）。边界更贴合车身轮廓，车窗和轮毂区域 mask 精度提升。}
  \end{columns}
\end{frame}

% ==================== 第 7 页：可视化案例二 ====================
\begin{frame}[t]{7. 可视化案例二：小型目标}
  \footnotesize
  \linespread{0.92}\selectfont

  \vspace{-0.2em}
  \centering
  \fcolorbox{black!25}{white}{%
    \includegraphics[width=0.80\linewidth,height=0.60\textheight,keepaspectratio]{pic/sam2_case2_compare.png}}

  \vspace{0.12em}
  \begin{columns}[T,onlytextwidth]
    \column{0.50\textwidth}
      \tcard{deepred}{1.10cm}{Baseline 问题}{%
        \small SAD=3843。微小目标（bbox 仅 117$\times$133）容易误选背景区域，mask 几乎全错。}
    \column{0.50\textwidth}
      \tcard{softgreen}{1.10cm}{优化后改善}{%
        \small SAD=348（$\downarrow$90.9\%）。Ensemble 重排序有效纠正了候选选择错误，精确捕获目标。}
  \end{columns}
\end{frame}

% ==================== 第 8 页：可视化案例三 ====================
\begin{frame}[t]{8. 可视化案例三：细结构物体}
  \footnotesize
  \linespread{0.92}\selectfont

  \vspace{-0.2em}
  \centering
  \fcolorbox{black!25}{white}{%
    \includegraphics[width=0.80\linewidth,height=0.60\textheight,keepaspectratio]{pic/sam2_case3_compare.png}}

  \vspace{0.12em}
  \begin{columns}[T,onlytextwidth]
    \column{0.50\textwidth}
      \tcard{deepred}{1.10cm}{Baseline 问题}{%
        \small SAD=2562。细结构处有 mask 碎片和空洞，边界不完整。}
    \column{0.50\textwidth}
      \tcard{softgreen}{1.10cm}{优化后改善}{%
        \small SAD=374（$\downarrow$85.4\%）。Mask 更完整，碎片显著减少，细边缘更清晰。}
  \end{columns}
\end{frame}

% ---- 本节结束 ----
 引入，
%    继承主文件导言区的全部宏 / 配色 / 字体，无需自带 \documentclass 或 \begin{document}。
%  · 内容：模型选型 → Prompt Ensemble + Rerank → 多尺度融合 + 消融 → 可视化
% ============================================================================
\section{上游分割}

% ==================== 第 1 页：模型选型与系统定位 ====================
\begin{frame}[t]{1. 上游模型选型：为什么用 SAM2？}
  \footnotesize
  \linespread{0.90}\selectfont

  \begin{columns}[T,onlytextwidth]
    \column{0.50\textwidth}
      \tcard{softorange}{2.00cm}{定位需求分析}{%
        \textbf{系统需求：}给定图像，输出 bbox / points，供 ZIM 做精细 matting。\\
        \textbf{关键约束：}
        \begin{itemize}
          \setlength{\itemsep}{0.02em}
          \item \hilite{zero-training} — 不微调模型参数；
          \item 推理高效 — 单卡秒级；
          \item bbox prompt 原生支持 — 与下游接口统一。
        \end{itemize}
        \vspace{0.02em}
        \textbf{不选 SAM3：}模型更大、推理慢约 2--3 倍，bbox-only 提升有限。\\
        \textbf{不选 Grounded-SAM-2：}额外引入 text encoder，增加显存开销。
      }
    \column{0.50\textwidth}
      \tcard{accentblue}{2.00cm}{SAM2 适配性}{%
        \begin{itemize}
          \setlength{\itemsep}{0.04em}
          \item \textbf{Hiera 架构} — 轻量高效，Hiera-Base+ 仅 80.8M 参数；
          \item \textbf{native bbox + points} — 无需适配层；
          \item \textbf{多候选 mask + IoU 评分} — 为 rerank 提供置信度；
          \item \textbf{实测性能} — A10 上约 1.16 it/s（bbox-only）；
          \item 社区活跃，checkpoint 可下载。
        \end{itemize}
      }
  \end{columns}

  \vspace{0.15em}
  \notestrip{accentblue}{\centering\scriptsize
  \textbf{定位：}SAM2 负责回答 \hilite{``目标在哪''} — 输出 bbox + points + 候选 mask；精细软 alpha 留给下游 ZIM + TTA-GF + BSP。}
\end{frame}

% ==================== 第 2 页：Prompt Ensemble + Mask Rerank ====================
\begin{frame}[t]{2. 核心优化一：Prompt Ensemble 与 IoU 引导重排序}
  \scriptsize
  \linespread{0.90}\selectfont

  \notestrip{deepred}{\tiny
  \textbf{动机：}SAM2 单次 bbox 推理易受局部纹理干扰，可能选到次优 mask。我们提出 \hilite{Prompt Ensemble + IoU 引导重排序}：多组 prompt 策略生成候选 mask，再利用 SAM2 内置 IoU 与结构先验选出最优。}

  \vspace{0.35em}
  \centering
  \resizebox{0.92\linewidth}{!}{%
  \begin{tikzpicture}[
    x=1cm, y=1cm,
    >={Latex[length=2.5mm,width=2mm]},
    every node/.style={font=\tiny},
    softsh/.style={blur shadow={shadow blur steps=5, shadow xshift=0.4pt,
                   shadow yshift=-0.5pt, shadow opacity=24, shadow blur radius=1.5pt}},
    nodebase/.style={softsh, rounded corners=3pt, line width=0.9pt, align=center, inner sep=3pt},
    io/.style={nodebase, draw=black!50, fill=black!5, minimum width=1.5cm, minimum height=0.58cm},
    sam/.style={nodebase, draw=accentblue!82!black, fill=accentblue!14, minimum width=1.55cm, minimum height=0.58cm},
    cand/.style={nodebase, draw=softorange!72!black, fill=softorange!12, minimum width=1.2cm, minimum height=0.58cm},
    ours/.style={nodebase, draw=nankai!88!black, fill=nankai!14, minimum width=1.35cm, minimum height=0.58cm},
    best/.style={nodebase, draw=softgreen!72!black, fill=softgreen!16, minimum width=1.35cm, minimum height=0.58cm,
                 line width=1.2pt},
    pipe/.style={-{Latex[length=2.4mm,width=2mm]}, line width=1pt, draw=nankai!80!black, rounded corners=2pt},
    tag/.style={font=\tiny\bfseries, text=nankai!88!black}
  ]
    % ---- top branch ----
    \node[io]   (img)  at (0,1.15) {\textbf{原图 + bbox}};
    \node[sam]  (sam0) at (2.2,1.15) {\textbf{SAM2}\\[0.5pt]{\tiny bbox-only}};
    \node[cand] (c0)   at (4.2,1.15) {\tiny 3 candidates};

    % ---- bottom branch ----
    \node[io]   (img2) at (0,0) {\textbf{原图 + bbox}\\[0.5pt]{\tiny + positive points}};
    \node[sam]  (sam1) at (2.2,0) {\textbf{SAM2}\\[0.5pt]{\tiny bbox + pos}};
    \node[cand] (c1)   at (4.2,0) {\tiny up to 3 candidates};

    % ---- merge into rerank ----
    \node[ours] (rerank) at (6.3,0.575) {\textbf{Ensemble}\\[0.5pt]{\tiny + Rerank}};
    \node[best] (out)    at (8.3,0.575) {\textbf{最优 mask}\\[0.5pt]{\tiny $\alpha$}};

    % ---- arrows ----
    \draw[pipe] (img)  -- (sam0);
    \draw[pipe] (sam0) -- (c0);
    \draw[pipe] (img2) -- (sam1);
    \draw[pipe] (sam1) -- (c1);
    \draw[pipe] (c0.east) -| (rerank.north);
    \draw[pipe] (c1.east) -| (rerank.south);
    \draw[pipe] (rerank) -- (out);

    % ---- grouping bands ----
    \begin{scope}[on background layer]
      \node[fit=(img)(sam0)(c0)(img2)(sam1)(c1), rounded corners=6pt,
            draw=accentblue!45, line width=0.8pt, fill=accentblue!5, inner sep=8pt] (grpA) {};
      \node[fit=(rerank)(out), rounded corners=6pt,
            draw=nankai!52, line width=0.8pt, fill=nankai!6, inner sep=8pt] (grpB) {};
    \end{scope}
    \node[tag, anchor=south west] at ([xshift=6pt,yshift=-1pt]grpA.north west)
      {多策略候选生成 · Multi-prompt generation};
    \node[tag, anchor=south east] at ([xshift=-6pt,yshift=-1pt]grpB.north east)
      {置信度重排序 · Confidence-aware rerank};
  \end{tikzpicture}}

  \vspace{0.04em}
  \begin{columns}[T,onlytextwidth]
    \column{0.24\textwidth}
      \ccard{nankai}{0.95cm}{%
      {\fontsize{6}{7}\selectfont\textbf{\textcolor{nankai!88!black}{重排序公式}}}\\[1pt]
      {\fontsize{5.5}{6.5}\selectfont $S(m)=0.55\cdot\text{IoU}(m)+0.25\cdot\text{conn}(m)+0.20\cdot\text{bbox\_cov}(m)$}}
    \column{0.25\textwidth}
      \ccard{accentblue}{0.95cm}{%
      {\fontsize{6}{7}\selectfont\textbf{\textcolor{accentblue!85!black}{SAM2 IoU (55\%)}}}\\[1pt]
      {\fontsize{5.5}{6.5}\selectfont 模型自身的 mask 质量置信度；最可靠信号}}
    \column{0.25\textwidth}
      \ccard{softorange}{0.95cm}{%
      {\fontsize{6}{7}\selectfont\textbf{\textcolor{softorange!75!black}{连通性 (25\%)}}}\\[1pt]
      {\fontsize{5.5}{6.5}\selectfont 连通分量数衡量 mask 碎片化；惩罚零散碎片}}
    \column{0.26\textwidth}
      \ccard{softgreen}{0.95cm}{%
      {\fontsize{6}{7}\selectfont\textbf{\textcolor{softgreen!60!black}{bbox 覆盖 (20\%)}}}\\[1pt]
      {\fontsize{5.5}{6.5}\selectfont mask 应填充 bbox 的 10--70\%；偏离的压低}}
  \end{columns}
\end{frame}

% ==================== 第 3 页：多尺度推理 ====================
\begin{frame}[t]{3. 核心优化二：多尺度推理}
  \scriptsize
  \linespread{0.92}\selectfont

  \notestrip{deepred}{\scriptsize
  \textbf{动机：}单尺度 SAM2 推理对图像分辨率敏感，边界处可能产生偏差。我们引入 \hilite{多尺度推理}：在多个尺度分别运行 SAM2，融合结果以降低单尺度偏差。}

  \vspace{0.70em}
  \begin{columns}[T,onlytextwidth]
    \column{0.50\textwidth}
      \tcard{accentblue}{3.20cm}{多尺度推理原理}{%
        \begin{itemize}
          \setlength{\itemsep}{0.08em}
          \item 在 3 个尺度 $\{0.75\times, 1.0\times, 1.25\times\}$ 分别运行 SAM2；
          \item 将各尺度 mask resize 回原图分辨率；
          \item 平均融合：$\alpha_{\text{fused}}=\frac{1}{3}\sum_{s}\alpha_s$；
          \item 单一尺度的边界偏差被多数投票稀释；
          \item 不增加模型参数，纯推理时后处理。
        \end{itemize}
        \vspace{0.10em}
        \textbf{与 Ensemble 叠加：}\hilite{每尺度各自 ensemble rerank}，再跨尺度平均 —— 两阶段互补，进一步提升鲁棒性。}
    \column{0.50\textwidth}
      \tcard{nankai}{3.20cm}{直观理解}{%
        \begin{itemize}
          \setlength{\itemsep}{0.10em}
          \item \textbf{放大（1.25×）：}捕捉更精细的边界细节；
          \item \textbf{原始（1.0×）：}保持模型训练分布下的最优精度；
          \item \textbf{缩小（0.75×）：}获得更大感受野，减少局部歧义。
        \end{itemize}
        \vspace{0.15em}
        三个尺度的 mask 经 resize + 平均后，\hilite{单尺度的随机偏差被平滑}，类似集成学习中 bagging 的思路。

        \vspace{0.12em}
        \textbf{代价：}推理时间约 $\times$3（3 个尺度），但与 ensemble 组合时收益大于开销。}
  \end{columns}
\end{frame}

% ==================== 第 4 页：消融实验设定 ====================
\begin{frame}[t]{4. 消融实验设定}
  \footnotesize
  \linespread{0.92}\selectfont

  \begin{columns}[T,onlytextwidth]
    \column{0.48\textwidth}
      \tcard{softgray}{2.60cm}{数据集与模型}{%
        \begin{itemize}
          \setlength{\itemsep}{0.08em}
          \item \textbf{数据集：}MicroMat-3K valid706 子集
          \begin{itemize}
            \item 706 张图像，64 张基础图像
            \item 每张配有 official bbox + GT alpha
          \end{itemize}
          \item \textbf{模型：}SAM2.1 Hiera-Base+（80.8M 参数）
        \end{itemize}
      }
      \vspace{0.15em}
      \tcard{accentblue}{1.90cm}{评测指标}{%
        \begin{itemize}
          \setlength{\itemsep}{0.06em}
          \item \textbf{SAD}（Sum of Absolute Differences）—— 绝对误差总和
          \item \textbf{MSE} —— 均方误差
          \item \textbf{MAE$\times$1000} —— 平均绝对误差（$\times$1000）
          \item \textbf{Boundary SAD} —— 边界带（5px）SAD，衡量边缘质量
        \end{itemize}
      }
    \column{0.52\textwidth}
      \tcard{nankai}{4.65cm}{消融方法设计（8 种组合）}{%
        \textbf{核心思路：}控制变量法，逐项检验各部分贡献。
        \vspace{0.08em}
        \begin{itemize}
          \setlength{\itemsep}{0.05em}
          \item \textbf{bbox\_binary} — 基线：bbox-only 单次推理 + 硬二值化
          \item \textbf{+ guided} — 单路 + guided filter 边缘平滑
          \item \textbf{+ multiscale} — 单路 + 3 尺度融合
          \item \textbf{+ multiscale\_guided} — 单路 + 多尺度 + guided
          \item \textbf{+ ensemble\_rerank} — 双 prompt 候选 + IoU 重排序
          \item \textbf{+ ensemble\_guided} — ensemble + guided
          \item \textbf{+ ensemble\_multiscale} — ensemble + 多尺度
          \item \textbf{+ ensemble\_multiscale\_guided} — 全部叠加
        \end{itemize}
        \vspace{0.08em}
        从 1 到 8 逐步叠加优化，分析各模块的\hilite{独立贡献与交互效应}。}
  \end{columns}
\end{frame}

% ==================== 第 5 页：消融结果 ====================
\begin{frame}[t]{5. 消融实验结果（valid706）}
  \scriptsize
  \linespread{0.92}\selectfont

  \begin{columns}[T,onlytextwidth]
    \column{0.58\textwidth}
      \begin{block}{完整消融对比表}
        \vspace{-0.3em}
        \centering
        {\fontsize{7}{8}\selectfont
        \setlength{\tabcolsep}{2.2pt}
        \renewcommand{\arraystretch}{1.04}
        \resizebox{0.99\linewidth}{!}{%
        \begin{tabular}{lrrrr}
          \toprule
          \textbf{方法} & \textbf{SAD $\downarrow$} & \textbf{MSE $\downarrow$} & \textbf{MAE$\times$1000 $\downarrow$} & \textbf{Bnd-SAD $\downarrow$} \\
          \midrule
          \textbf{ensemble\_guided (最终)} & \textbf{2774.0} & \textbf{0.000747} & \textbf{0.981} & \textbf{1215.5} \\
          ensemble\_multiscale\_guided & 2915.4 & 0.000742 & 1.033 & 1188.8 \\
          ensemble\_rerank & 2930.9 & 0.000843 & 1.037 & 1380.1 \\
          ensemble\_multiscale & 2997.4 & 0.000796 & 1.062 & 1279.2 \\
          multiscale\_guided & 3145.3 & 0.000828 & 1.089 & 1249.0 \\
          guided & 3161.9 & 0.000845 & 1.095 & 1271.2 \\
          multiscale & 3228.3 & 0.000894 & 1.119 & 1340.7 \\
          \textbf{bbox\_binary (基线)} & \textbf{3313.5} & \textbf{0.000956} & \textbf{1.149} & \textbf{1431.5} \\
          \bottomrule
        \end{tabular}}}
      \end{block}
    \column{0.42\textwidth}
      \tcard{softgreen}{2.05cm}{关键发现}{%
        \begin{itemize}
          \setlength{\itemsep}{0.08em}
          \item \hilite{Ensemble rerank 单点最有效}（SAD$\downarrow$11.5\% vs 基线）
          \item Guided filter 对 Boundary 改善显著（1432$\to$1271）
          \item Multiscale 单独增益有限（3313$\to$3228）
          \item 三者叠加达最优：SAD$\downarrow$16.3\%，MAE$\downarrow$14.6\%
        \end{itemize}
      }

      \vspace{0.25em}
      \tcard{nankai}{1.30cm}{最终提交方法}{%
        \centering
        Ensemble Rerank + Guided Filter}

      \vspace{0.6em}
      \begin{columns}[T,onlytextwidth]
        \column{0.33\textwidth}
          \metricchip{softgreen}{SAD 改善{\tiny\ (vs 基线)}}{$-16.3\%$}
        \column{0.33\textwidth}
          \metricchip{softgreen}{MAE 改善{\tiny\ (vs 基线)}}{$-14.6\%$}
        \column{0.34\textwidth}
          \metricchip{nankai}{Boundary 改善{\tiny\ (vs 基线)}}{$-15.1\%$}
      \end{columns}
  \end{columns}
\end{frame}

% ==================== 第 6 页：可视化案例一 ====================
\begin{frame}[t]{6. 可视化案例一：大型目标（汽车）}
  \footnotesize
  \linespread{0.92}\selectfont

  \vspace{-0.2em}
  \centering
  \fcolorbox{black!25}{white}{%
    \includegraphics[width=0.92\linewidth,height=0.60\textheight,keepaspectratio]{pic/sam2_case1_compare.png}}

  \vspace{0.12em}
  \begin{columns}[T,onlytextwidth]
    \column{0.50\textwidth}
      \tcard{deepred}{1.10cm}{Baseline 问题}{%
        \small SAD=13559。车轮和车窗边缘有明显的 mask 溢出，导致前景区域偏大。}
    \column{0.50\textwidth}
      \tcard{softgreen}{1.10cm}{优化后改善}{%
        \small SAD=9010（$\downarrow$33.5\%）。边界更贴合车身轮廓，车窗和轮毂区域 mask 精度提升。}
  \end{columns}
\end{frame}

% ==================== 第 7 页：可视化案例二 ====================
\begin{frame}[t]{7. 可视化案例二：小型目标}
  \footnotesize
  \linespread{0.92}\selectfont

  \vspace{-0.2em}
  \centering
  \fcolorbox{black!25}{white}{%
    \includegraphics[width=0.80\linewidth,height=0.60\textheight,keepaspectratio]{pic/sam2_case2_compare.png}}

  \vspace{0.12em}
  \begin{columns}[T,onlytextwidth]
    \column{0.50\textwidth}
      \tcard{deepred}{1.10cm}{Baseline 问题}{%
        \small SAD=3843。微小目标（bbox 仅 117$\times$133）容易误选背景区域，mask 几乎全错。}
    \column{0.50\textwidth}
      \tcard{softgreen}{1.10cm}{优化后改善}{%
        \small SAD=348（$\downarrow$90.9\%）。Ensemble 重排序有效纠正了候选选择错误，精确捕获目标。}
  \end{columns}
\end{frame}

% ==================== 第 8 页：可视化案例三 ====================
\begin{frame}[t]{8. 可视化案例三：细结构物体}
  \footnotesize
  \linespread{0.92}\selectfont

  \vspace{-0.2em}
  \centering
  \fcolorbox{black!25}{white}{%
    \includegraphics[width=0.80\linewidth,height=0.60\textheight,keepaspectratio]{pic/sam2_case3_compare.png}}

  \vspace{0.12em}
  \begin{columns}[T,onlytextwidth]
    \column{0.50\textwidth}
      \tcard{deepred}{1.10cm}{Baseline 问题}{%
        \small SAD=2562。细结构处有 mask 碎片和空洞，边界不完整。}
    \column{0.50\textwidth}
      \tcard{softgreen}{1.10cm}{优化后改善}{%
        \small SAD=374（$\downarrow$85.4\%）。Mask 更完整，碎片显著减少，细边缘更清晰。}
  \end{columns}
\end{frame}

% ---- 本节结束 ----
